% !PRINTED 163
\weberpage{163}
%<<CONT
of unity, and $e$ shall be called the \gesperrt{index of the invariant} $\mathit\Psi(x)$.
If $\varepsilon$ is a primitive $e$th root of unity, we can set
\[
\alpha=\varepsilon^a,\quad \beta=\varepsilon^b,\quad \gamma=\varepsilon^c \ldots
\]
and the exponents $a,b,c\ldots$ can have no common
divisor with $e$. From this it follows that the
integers $x,y,z,\ldots$ can be determined so that
\[
ax+by+cz+\cdots \equiv 1\ (\mathrm{mod}\ e)
\]
holds (Bd. I, \S.\,118). Now since, on account of the group nature, among
the factors $\alpha,\beta,\gamma,\ldots$ the number
\[
\alpha^x\,\beta^y\,\gamma^z \cdots = \varepsilon
\]
also occurs, it follows that the totality of the factors $\alpha,\beta,\gamma\ldots$
must coincide with the powers of $\varepsilon$:
\[
1,\ \varepsilon,\ \varepsilon^2,\ \ldots\ \varepsilon^{e-1}
\]

If we seek in $S$ all substitutions $A$ that satisfy the condition
\weq{4}{\mathit\Psi[A(x)] = \mathit\Psi(x)}
to which the identical substitution certainly belongs, we
obtain a new group $T$, which is a divisor of $S$.
If $E$ then denotes a substitution of $S$ satisfying the condition
\weq{5}{\mathit\Psi[E(x)] = \varepsilon\,\mathit\Psi(x)}
then all substitutions $AE$
and $EA$ have the same property, and we obtain the decomposition
of $S$ into the subsidiary groups:
\weq{6}{S = T + TE + TE^2 + \cdots + TE^{e-1};}
the index of the divisor $T$ is therefore $=\varepsilon$. % [sic]
At the same time it follows further, since every substitution $E^{-1}AE$ satisfies condition (4), that
\weq{7}{E^{-1}TE = T,}
from which it appears that $T$ is a normal divisor of $S$. We
have thus proved the theorem:

3. \gesperrt{If $\mathit\Psi(x)$ is an invariant of the group $S$ of
index $e$, then all substitutions of $S$
by which $\mathit\Psi(x)$ remains unchanged form a normal divisor
 $T$ of $S$ of index $e$.}

For the group $T$, $\mathit\Psi(x)$ is an absolute invariant. The invariants
 of index 1 are the absolute invariants of $S$.
The group $T$ is called the \gesperrt{group belonging to the invariant
$\mathit\Psi(x)$.}
